\documentclass{article}
\input{../../lib.sty}

\title{\texttt{fn permute\_and\_flip}}
\author{\MichaelShoemateAuthor}

\begin{document}
\maketitle\contrib
This document proves soundness of \rustdoc{measurements/noisy_top_k}{permute\_and\_flip} \cite{mckenna2020permute} in \asOfCommit{mod.rs}{7d1a20076}.
\texttt{permute\_and\_flip} noisily selects the index of the greatest score from a vector of input scores.

\section{Hoare Triple}
\subsection*{Precondition}
\subsubsection*{Compiler-verified}
\textit{Types consistent with pseudocode.}

\subsubsection*{Caller-verified}
\texttt{x} is nonempty and \texttt{scale} is nonnegative.

\subsection*{Pseudocode}
\label{sec:python-pseudocode}
\lstinputlisting[language=Python,firstline=2,escapechar=`]{./pseudocode/permute_and_flip.py}

\subsection*{Postcondition}
\begin{theorem}
    \label{postcondition}
    \begin{itemize}
        \item If \texttt{scale} is zero, returns an index of a greatest score, breaking ties deterministically.
        \item Otherwise, if replacement is set, returns a sample from $\mathcal{M}_{EM}$ (as defined in MS2020 Definition 4),
            otherwise returns a sample from $\mathcal{M}_{PF}$ (as defined in MS2020 Lemma 1),
            where $\texttt{scale} = \frac{2 \cdot \Delta}{\epsilon}$.
        \item Errors are data-independent, except for exhaustion of entropy.
    \end{itemize}
\end{theorem}

\begin{proof}
    If \texttt{scale} is zero, the function returns a maximizer directly.
    Otherwise, when replacement is not set, swapping elements on line~\ref{shuffle}
    applies an online Fisher-Yates shuffle up to and including index \texttt{left}.

    Substituting $\texttt{scale} = \frac{2\Delta}{\epsilon}$,
    the argument to \rustdoc{traits/samplers/cks20/fn}{sample\_bernoulli\_exp} is then $\frac{\epsilon}{2\Delta}(q_* - q_r)$,
    which is non-negative, satisfying the precondition of \rustdoc{traits/samplers/cks20/fn}{sample\_bernoulli\_exp}.
    Therefore by the postcondition of \rustdoc{traits/samplers/cks20/fn}{sample\_bernoulli\_exp},
    the response is a sample from $\mathrm{Bern}(\exp(-x))$, where $x = \frac{\epsilon}{2\Delta}(q_* - q_r)$.
    Therefore the response is a sample from $\mathrm{Bern}(\exp(\frac{\epsilon}{2\Delta}(q_r - q_*)))$,
    which, without replacement, is equivalent to Algorithm~1 in \cite{mckenna2020permute}.
    At least one candidate has score $q_*$ and is accepted with probability one.
    With replacement, \texttt{left} is always zero.
    Conditional on any previous trials, the candidate array is a permutation,
    so swapping position zero with a fresh uniform index proposes each candidate with probability $1/n$, where $n=|x|$.
    The subsequent Bernoulli trial uses fresh independent randomness.
    The probability of accepting candidate $r$ in one trial is
    $p_r=\exp((q_r-q_*)/\texttt{scale})/n$.
    The total acceptance probability $p=\sum_r p_r$ is at least $1/n$,
    since a maximizer is accepted with probability one.
    Thus the loop terminates almost surely, after $1/p\leq n$ trials in expectation,
    and the probability of returning candidate $r$ is
    $\sum_{j=0}^{\infty}(1-p)^j p_r=p_r/p$,
    as required for the exponential mechanism.

    The only source of error is due to entropy exhaustion.
\end{proof}

\bibliographystyle{plain}
\bibliography{references.bib}

\end{document}
